\begin{lemma}[At most one stored color survives per series round]
For every certified Tashkinov series and every data index \(i\),
\[
 |D_i|\leq i.
\]
Here \(D_i\) is the finite set \noderef{TashkinovStoredColors}; in
particular, \(D_0=\varnothing\).  The natural number on the right is the
value of the zero-based Lean index.
\end{lemma}
\begin{proof}
We prove the assertion by induction on the natural value of \(i\).  The
definition \noderef{TashkinovStoredColors} gives \(D_0=\varnothing\), so
the bound is equality at zero.

For the induction step write \(i=n\), with \(n\geq1\), and write the
preceding zero-based round as \(a=n-1\).  Thus the current positive data
slot is \(a.\mathrm{succ}\).  We first prove the finite-set recurrence
\[
 D_n\subseteq D_{n-1}\cup\{x_n\}                               \tag{1}
\]
for some color \(x_n\), allowing the right-hand singleton to be omitted
in the RE case.

Use the certified step supplied by \noderef{TashkinovSeries} at this
positive slot, and distinguish the constructors in
\noderef{TashkinovSeriesStep}.

In an RE, the current coloring equals the preceding coloring.  Its two
stored colors are the gamma and delta of an explicitly recorded earlier
PE slot.  Consequently the current stored pair is already contained in
the union of earlier stored sets; it adds no color to that union.  The
current tree contains the preceding tree as an ordered initial segment,
by \noderef{TreeSequenceInitialSegment}, so under the unchanged coloring
its tree-missing set contains the preceding tree-missing set.  Subtracting
this larger set from the same accumulated stored union gives
\[
 D_n\subseteq D_{n-1}.                                         \tag{2}
\]
This is the exact recurrence supplied by the witness-carrying RE
constructor.  It is also the content behind the paper's notation that
an RE reuses its active stored pair.

In an SE, let \(\delta_n\) be the selected defective color.  The
\noderef{MaximumDefectiveChoice} makes the current coloring stable
relative to the preceding coloring on \(T_n\).  The vertex clause of
\noderef{StableColoring} therefore leaves the missing-color union on
\(T_n\) unchanged.  The only new stored set is
\(S_n=\{\delta_n\}\).  Earlier tree vertices remain in \(T_n\), again by
\noderef{TreeSequenceInitialSegment}; hence, after expanding the two set
differences in \noderef{TashkinovStoredColors},
\[
 D_n\subseteq D_{n-1}\cup\{\delta_n\}.                          \tag{3}
\]

It remains to consider a PE.  Let \(\pi\) be its selected coloring,
\(\rho\) its second stable coloring, and let its stored pair be
\(S_n=\{\delta_n,\gamma_n\}\).  The two applications of
\noderef{StableColoring} identify the missing-color set of \(T_n\)
under the preceding coloring, \(\pi\), and \(\rho\).  The rule
\noderef{ParallelExtensionColor} says that \(\gamma_n\) is missing at
the supporting vertex under \(\rho\), while the controlled connecting
edge has color \(\delta_n\) there.  The PE's \noderef{TreeExit} witness,
with the conditions in \noderef{TreeExitPath}, is a simple path from
outside \(T_n\) to that support and has no other tree vertex.  Both path
endpoints have degree one in the two-color incidence graph, and every
degree is at most two by
\noderef{TwoColorIncidenceDegreeAtMostTwo}.  A shortest-attachment
argument shows that this path exhausts the component: an extra branch
would give degree three internally or degree two at an endpoint.  Hence
the switched component meets \(T_n\) only at the supporting vertex.
The definition \noderef{TwoColorComponentSwap} of the current coloring
therefore replaces the missing color \(\gamma_n\) by \(\delta_n\) at
that vertex and loses no other tree-missing color.  In particular,
\[
 \bigl(\overline\varphi_{n-1}(T_n)\setminus\{\gamma_n\}\bigr)
       \cup\{\delta_n\}
 \subseteq \overline\varphi_n(T_n).                             \tag{4}
\]
This is the PE instance of the first propagation established in
\noderef{HkuMissingStoredMonotone}.

Let \(U_{n-1}=\bigcup_{h\leq n-1}S_h\).  Using (4) and finite-set
algebra,
\[
\begin{aligned}
 D_n
 &=\bigl(U_{n-1}\cup\{\delta_n,\gamma_n\}\bigr)
      \setminus\overline\varphi_n(T_n)\\
 &\subseteq
   \bigl(U_{n-1}\setminus\overline\varphi_{n-1}(T_n)\bigr)
      \cup\{\gamma_n\}\\
 &\subseteq D_{n-1}\cup\{\gamma_n\}.
\end{aligned}                                                    \tag{5}
\]
For the last inclusion, the preceding tree is an initial segment of
\(T_n\), so a color missing on the preceding tree is also missing on
\(T_n\).  This completes all cases of (1).

For finite sets, inclusion (1), monotonicity of cardinality, and
\(|A\cup\{x\}|\leq |A|+1\) give
\[
 |D_n|\leq |D_{n-1}|+1.
\]
The induction hypothesis gives \(|D_{n-1}|\leq n-1\), and hence
\(|D_n|\leq n\).  This proves the bound for every allowed finite index.
\end{proof}
